\chapter{Classical Inference}
\label{cap:inferencia}
% ── index ───────────────────────────────────────────────────────
\index{t-test|textbf}
\index{ttest@\texttt{ttest()}|textbf}
\index{ANOVA|textbf}
\index{anova@\texttt{anova()}|textbf}
\index{Kruskal-Wallis, test of|textbf}
\index{kruskal_wallis@\texttt{kruskal\_wallis()}|textbf}
\index{Welch, test of|textbf}
\index{nonparametric, test}

% ─────────────────────────────────────────────────────────────────────────────
\section{Parametric and non-parametric hypothesis tests}

The tests presented in this chapter answer the question: \emph{are the data
consistent with a hypothesis about the population distribution?}

\subsection*{Student's t-test}

For a sample of size $n$ with unknown variance, the statistic:
\[
  t = \frac{\bar{x} - \mu_0}{s / \sqrt{n}} \sim t_{n-1}
\]
tests $H_0: \mu = \mu_0$. For two independent samples, Welch's test
uses approximate degrees of freedom by the Satterthwaite method, without
assuming homoscedasticity.

\subsection*{One-way ANOVA}

Generalizes the $t$-test to $K > 2$ groups. Decomposes total variation:
\[
  \text{SS}_{\text{total}} = \text{SS}_{\text{between}} + \text{SS}_{\text{within}}
\]
\[
  F = \frac{\text{SS}_{\text{between}} / (K-1)}{\text{SS}_{\text{within}} / (N-K)}
  \sim F_{K-1,\, N-K}
\]
$H_0$: all group means are equal.

\subsection*{Kruskal-Wallis}

Non-parametric alternative to one-way ANOVA. Replaces observations with
\emph{ranks} and compares rank sums across groups. Statistic:
\[
  H = \frac{12}{N(N+1)} \sum_{k=1}^K \frac{R_k^2}{n_k} - 3(N+1) \sim \chi^2_{K-1}
\]
valid even without normality, as long as the distributions have a similar shape.

% ─────────────────────────────────────────────────────────────────────────────
\section{DGP Verification}

DGP with 3 groups of different sizes; the expected value increases linearly
with the group. $N=300$, $\sigma = 1$.

\begin{lstlisting}[caption={Hypothesis tests}]
seed(42)
...
generate df g   = (_n <= 150) * 1 + (_n > 150) * 2
generate df yg  = 2.0 + e + (g - 1) * 0.5  // g=1: mu=2.0, g=2: mu=2.5

// One-sample t-test: H0: mu = 2.0
ttest(df, y, mu=2.0)

// Two-sample t-test (Welch): H0: mu(g1) = mu(g2)
ttest(df, yg, by=g)

// Paired: before vs after
ttest(df, y_antes, y_depois, paired=true)

// ANOVA: H0: mu1 = mu2 = mu3
anova(df, y3, by=k)

// Kruskal-Wallis (non-parametric)
kruskal_wallis(df, y3, by=k)
\end{lstlisting}

\subsection*{t-test — one sample}

\begin{lstlisting}[language={}, basicstyle=\ttfamily\small, frame=none,
  numbers=none, backgroundcolor=\color{codbg}]
One-sample t-test: y   H0: mean = 2.0
t = -0.6946   df = 299   p = 0.4879
\end{lstlisting}

Does not reject $H_0$ — sample mean 1{,}98 is compatible with 2{,}0.

\subsection*{t-test — two groups (Welch)}

\begin{lstlisting}[language={}, basicstyle=\ttfamily\small, frame=none,
  numbers=none, backgroundcolor=\color{codbg}]
Two-sample t-test (Welch): yg by g
  group 1: n=150  mean=1.971  SE=0.040
  group 2: n=150  mean=2.489  SE=0.041
Welch's t = -9.0552   df = 297.38   p = 0.0000
\end{lstlisting}

Rejects $H_0$: difference of 0{,}52 $\approx$ true difference of 0{,}50.

\subsection*{ANOVA}

\begin{lstlisting}[language={}, basicstyle=\ttfamily\small, frame=none,
  numbers=none, backgroundcolor=\color{codbg}]
=========================== One-Way ANOVA ============================
Source           SS     df      MS         F       P>F
----------------------------------------------------------------------
Between       51.195     2   25.597   104.037   0.0000
Within        73.074   297    0.246
Total        124.269   299
\end{lstlisting}

$F = 104{,}0$, $p < 0{,}001$: rejects $H_0$ decisively.

\subsection*{Kruskal-Wallis}

\begin{lstlisting}[language={}, basicstyle=\ttfamily\small, frame=none,
  numbers=none, backgroundcolor=\color{codbg}]
H = 120.38   df = 2   p = 0.0000
\end{lstlisting}

Confirms the ANOVA without assuming normality.

% ─────────────────────────────────────────────────────────────────────────────
\section{Syntax}

\begin{lstlisting}
ttest(df, var, mu=0.0)                  // one sample
ttest(df, var, by=group)                 // two groups (Welch)
ttest(df, var, by=group, unequal=false) // two groups (Student)
ttest(df, var1, var2, paired=true)      // paired
anova(df, var, by=group)                // one-way ANOVA
kruskal_wallis(df, var, by=group)       // Kruskal-Wallis
\end{lstlisting}

\begin{nota}
  Welch's test does not assume homoscedasticity and is preferable to the
  classical Student's $t$-test in almost all practical cases. The loss of power
  is negligible when variances are equal; the gain when they differ
  is substantial. If you wish to run the classical Student's t-test, use the
  option \texttt{unequal=false}.
\end{nota}
